\documentclass[12pt,a4paper]{article}

% -------------------------------------------------
% Basic packages
% -------------------------------------------------
\usepackage[a4paper,margin=1in]{geometry}
\usepackage{fontspec}
\usepackage{ctex}
\usepackage{amsmath,amssymb}
\usepackage{graphicx}
\usepackage{booktabs}
\usepackage{array}
\usepackage{longtable}
\usepackage{caption}
\usepackage{float}
\usepackage{setspace}
\usepackage{fancyhdr}
\usepackage{hyperref}
\usepackage{xcolor}
\usepackage{enumitem}
\usepackage{listings}

% -------------------------------------------------
% Table of contents setup
% -------------------------------------------------
\setcounter{tocdepth}{2}

% -------------------------------------------------
% Font setup
% -------------------------------------------------
\setmainfont{Noto Serif}
\setmonofont{DejaVu Sans Mono}

% -------------------------------------------------
% Line spacing and paragraph formatting
% -------------------------------------------------
\onehalfspacing
\setlength{\parindent}{2em}
\setlength{\parskip}{0.4em}

% -------------------------------------------------
% Header / footer
% -------------------------------------------------
\pagestyle{fancy}
\fancyhf{}
\fancyfoot[C]{\thepage}
\renewcommand{\headrulewidth}{0pt}

% -------------------------------------------------
% Hyperref setup
% -------------------------------------------------
\hypersetup{
    colorlinks=true,
    linkcolor=blue,
    urlcolor=blue,
    citecolor=blue
}

% -------------------------------------------------
% Code listing setup
% -------------------------------------------------
\lstset{
    basicstyle=\ttfamily\small,
    breaklines=true,
    frame=single,
    columns=fullflexible,
    showstringspaces=false
}

% -------------------------------------------------
% User-defined information
% -------------------------------------------------
\newcommand{\HomeworkTitle}{Computational Physics Ex\_3}
\newcommand{\StudentID}{2024141230044}
\newcommand{\StudentName}{范奕}

\begin{document}

% =================================================
% Title block
% =================================================
\begin{center}
    {\LARGE \textbf{\HomeworkTitle}}\\[1.2em]
    {\large Student ID: \StudentID}\\[0.5em]
    {\large Name: \StudentName}
\end{center}

\vspace{0.5em}
\noindent\rule{\textwidth}{0.8pt}
\vspace{0.5em}

\section*{Problem 2}

使用 Machin 公式
\[
\frac{\pi}{4}=4\arctan\left(\frac{1}{5}\right)-\arctan\left(\frac{1}{239}\right)
\]
并结合反正切级数
\[
\arctan(x)=x-\frac{x^3}{3}+\frac{x^5}{5}-\frac{x^7}{7}+\cdots
\]
计算 $\pi$ 的值，并与程序中的 $\pi$ 真值进行比较。

\newpage
\tableofcontents
\newpage

% =================================================
% Part 1
% =================================================
\section{问题说明}

Machin 公式是一种经典的高精度计算 $\pi$ 的方法。其优点在于：
\begin{enumerate}[label=\arabic*)]
    \item 将 $\pi$ 的计算转化为两个较小自变量的 $\arctan$ 求值；
    \item 当 $x=\frac{1}{5}$ 与 $x=\frac{1}{239}$ 时，级数收敛非常快；
    \item 配合 \texttt{quadmath.h} 中的 \texttt{\_\_float128} 高精度浮点类型，可以较稳定地得到 30 位小数结果。
\end{enumerate}

本题中采用如下思路：
\begin{enumerate}[label=\arabic*)]
    \item 用反正切幂级数计算 $\arctan(1/5)$；
    \item 用反正切幂级数计算 $\arctan(1/239)$；
    \item 代入 Machin 公式计算 $\pi$；
    \item 使用 \texttt{acosq(-1.0Q)} 作为程序中的 $\pi$ 真值；
    \item 比较两者的绝对误差。
\end{enumerate}

\section{算法原理}

\subsection{Machin 公式}

\[
\frac{\pi}{4}=4\arctan\left(\frac{1}{5}\right)-\arctan\left(\frac{1}{239}\right)
\]

因此可得
\[
\pi = 4\left(4\arctan\left(\frac{1}{5}\right)-\arctan\left(\frac{1}{239}\right)\right)
\]

\subsection{反正切级数展开}

当 $|x|\le 1$ 时，反正切函数有幂级数展开
\[
\arctan(x)=\sum_{n=0}^{\infty}(-1)^n\frac{x^{2n+1}}{2n+1}
\]

即
\[
\arctan(x)=x-\frac{x^3}{3}+\frac{x^5}{5}-\frac{x^7}{7}+\cdots
\]

由于本题中使用的 $x=\frac{1}{5}$ 与 $x=\frac{1}{239}$ 都很小，所以该级数收敛很快。程序中采用“当当前项绝对值小于 $10^{-40}$ 时停止”的策略，保证结果足够精确。

\section{程序代码}

\begin{lstlisting}[language=C]
#include <stdio.h>
#include <quadmath.h>

// 用反正切级数计算 arctan(x)
__float128 arctan_series(__float128 x)
{
    __float128 sum = 0.0Q;
    __float128 term = x;
    __float128 x2 = x * x;
    int n = 0;

    while (fabsq(term) > 1e-40Q) {
        sum += term / (2 * n + 1);
        term *= -x2;
        n++;
    }

    return sum;
}

int main()
{
    // Machin公式计算pi
    __float128 a = arctan_series(1.0Q / 5.0Q);
    __float128 b = arctan_series(1.0Q / 239.0Q);
    __float128 pi_calc = 4.0Q * (4.0Q * a - b);

    // pi真值：利用 acos(-1)
    __float128 pi_true = acosq(-1.0Q);

    // 绝对误差
    __float128 abs_error = fabsq(pi_calc - pi_true);

    char buf1[128], buf2[128], buf3[128];

    quadmath_snprintf(buf1, sizeof(buf1), "%.30Qf", pi_calc);
    quadmath_snprintf(buf2, sizeof(buf2), "%.30Qf", pi_true);
    quadmath_snprintf(buf3, sizeof(buf3), "%.40Qe", abs_error);

    printf("Machin公式计算得到的 pi = %s\n", buf1);
    printf("pi 真值(acos(-1))     = %s\n", buf2);
    printf("绝对误差              = %s\n", buf3);

    return 0;
}
\end{lstlisting}

\section{运行结果}

根据程序实际输出，结果如下：

\begin{lstlisting}
Machin公式计算得到的 pi = 3.141592653589793238462643383280
pi 真值(acos(-1))     = 3.141592653589793238462643383280
绝对误差              = 0.0000000000000000000000000000000000000000e+00
\end{lstlisting}

\section{结果分析}

\subsection{与真值对比}

从输出结果可见，利用 Machin 公式和反正切级数计算得到的 $\pi$ 值为
\[
3.141592653589793238462643383280
\]

而程序中通过 \texttt{acosq(-1.0Q)} 得到的 $\pi$ 真值也是
\[
3.141592653589793238462643383280
\]

二者在 30 位小数范围内完全一致，程序输出的绝对误差为
\[
0.0000000000000000000000000000000000000000\times 10^0
\]

也就是在当前输出精度下，误差为 0。

\subsection{误差来源分析}

虽然程序输出显示误差为 0，但这并不表示数学意义上误差绝对为零，而是说明：
\begin{enumerate}[label=\arabic*)]
    \item 在 \texttt{\_\_float128} 的表示精度范围内，两者已经一致；
    \item 输出格式只显示了有限位数，因此更小的误差无法显示出来；
    \item Machin 公式本身收敛快，数值稳定性较好，因此在本题精度要求下非常适合计算 $\pi$。
\end{enumerate}

\subsection{为什么 Machin 公式效果好}

Machin 公式的优势在于：它避免了直接对 $\pi$ 进行慢收敛级数计算，而是将问题拆成两个小参数的反正切计算。由于：
\[
\left|\frac{1}{239}\right| \ll 1
\]
所以 $\arctan(1/239)$ 的级数项衰减极快；同时 $\arctan(1/5)$ 也具有较快收敛速度。因此整体算法效率较高，且容易达到高精度。

\section{编译与运行说明}

在支持 \texttt{quadmath} 的 GNU 编译环境下，可使用如下命令编译：

\begin{lstlisting}[language=bash]
gcc machin_pi.c -o machin_pi -lquadmath
\end{lstlisting}

运行程序：

\begin{lstlisting}[language=bash]
./machin_pi
\end{lstlisting}

\section{结论}

本实验使用 Machin 公式
\[
\frac{\pi}{4}=4\arctan\left(\frac{1}{5}\right)-\arctan\left(\frac{1}{239}\right)
\]
以及反正切幂级数展开，实现了对 $\pi$ 的高精度计算。程序结果表明：

\begin{enumerate}[label=\arabic*)]
    \item 计算得到的 $\pi$ 值为 $3.141592653589793238462643383280$；
    \item 与程序真值 \texttt{acosq(-1.0Q)} 在 30 位小数范围内完全一致；
    \item 输出绝对误差为 0（在当前显示精度下）；
    \item Machin 公式收敛速度快、数值稳定，适合用于高精度计算 $\pi$。
\end{enumerate}

因此，本题所用算法能够满足“计算 $\pi$ 到 30 位小数”的要求，且结果正确可靠。

\end{document}
